\documentclass[11pt,a4paper]{article}
\usepackage[margin=2.5cm]{geometry}
\usepackage{amsmath,amssymb}
\usepackage{booktabs}
\usepackage{xcolor}
\usepackage{hyperref}
\usepackage{enumitem}
\usepackage{parskip}
\usepackage{titlesec}
\usepackage{mdframed}

% Colour scheme
\definecolor{reviewerblue}{RGB}{0,70,127}
\definecolor{responsegreen}{RGB}{0,100,50}
\definecolor{changebrown}{RGB}{139,69,19}
\definecolor{lightgray}{RGB}{245,245,245}

% Reviewer comment box
\newmdenv[
  backgroundcolor=lightgray,
  linecolor=reviewerblue,
  linewidth=1pt,
  leftmargin=0pt,
  rightmargin=0pt,
  innerleftmargin=10pt,
  innerrightmargin=10pt,
  innertopmargin=6pt,
  innerbottommargin=6pt,
  skipabove=8pt,
  skipbelow=4pt
]{reviewbox}

% Section formatting
\titleformat{\section}{\large\bfseries\color{reviewerblue}}{}{0em}{}[\titlerule]
\titleformat{\subsection}{\normalsize\bfseries\color{reviewerblue}}{}{0em}{}

% Macros
\newcommand{\RC}[1]{\begin{reviewbox}\textit{\textbf{Reviewer Comment:} #1}\end{reviewbox}}
\newcommand{\RS}[1]{\noindent\textcolor{responsegreen}{\textbf{Response:}} #1\par\medskip}
\newcommand{\CH}[1]{\noindent\textcolor{changebrown}{\textbf{Change in manuscript:}} #1\par\bigskip}

\begin{document}

% -----------------------------------------------------------------------
%  HEADER
% -----------------------------------------------------------------------
\begin{center}
  {\LARGE\bfseries Response to Reviewers}\\[1em]
  {\large Manuscript: \textit{Agentic SABRE: An Uncertainty-Aware Neuro-Symbolic
  Multi-Agent Framework for Adaptive Ransomware Detection}}\\[0.5em]
  {\normalsize Expert Systems With Applications --- Major Revision}\\[0.5em]
  {\normalsize \today}
\end{center}

\hrule
\bigskip

% -----------------------------------------------------------------------
%  GENERAL PREAMBLE
% -----------------------------------------------------------------------
We thank the Associate Editor and all four reviewers for their thorough and
constructive evaluations. The comments have led to substantive improvements in
the manuscript's clarity, technical rigour, and positioning. We have addressed
every point raised. The changes are highlighted in \textbf{bold blue} in the
revised manuscript. For conciseness, we state in each response exactly where in
the revised paper the change appears.

A brief summary of the main additions:
\begin{itemize}[leftmargin=1.5em,itemsep=2pt]
  \item A new ``\textit{What Makes This Different from Ensemble-Based Approaches}''
        subsection in the Introduction, distinguishing SABRE's cooperative agentic
        design from conventional ensemble voting.
  \item A 7-column comparison table in Related Work summarising how prior
        systems differ from SABRE across key design dimensions.
  \item Full variable definitions and a notation glossary before the first score-vector
        equation; elimination of a duplicate $\sigma_{\max}$ formula; relocation of the
        triage-objective $J(\tau,\kappa)$ to the Policy Analysis section to remove a
        forward reference spanning 15 equations.
  \item Explicit threshold-ordering constraints ($\tau \le \tau_{\text{high}}$,
        $\kappa_{\text{low}} \le \kappa$) added as a numbered equation.
  \item New discussion paragraphs on: AUC$=1.0$ generalisation risk; the
        accuracy/AUC discrepancy of the Behavioural Agent; threshold sensitivity
        for $\tau = 0.999999$; cross-source fusion rationale; CTGAN label
        clarification; Monte Carlo Dropout vs.\ deep ensembles; real-world
        executable constraints for counterfactual costs; and a runtime/complexity
        subsection.
  \item Section-opening purpose sentences in all 13 sections.
\end{itemize}

\bigskip
\hrule
\bigskip

% -----------------------------------------------------------------------
%  ASSOCIATE EDITOR
% -----------------------------------------------------------------------
\section{Associate Editor}

\RC{The authors should explicitly mention the significant contributions of the
manuscript. The novelty of the paper is not highlighted.}

\RS{We agree that the novelty was insufficiently foregrounded. We have added a new
subsection titled \textit{``What Makes This Different from Ensemble-Based Approaches''}
immediately after the Contributions list in Section~I. This subsection explains,
in concrete technical terms, the three properties that distinguish Agentic SABRE
from conventional ensemble detectors: (1) structurally heterogeneous input spaces
($\mathbb{R}^{384}$ semantic embeddings versus $\mathbb{R}^{16}$ I/O telemetry),
so agents are complementary by design rather than redundant; (2) the triage decision
conditions on epistemic \emph{uncertainty}, not only on posterior probabilities;
and (3) the triage policy is an \emph{executable control mechanism}---not soft label
aggregation---that can formally withhold autonomous action under ambiguity.}

\CH{New subsection ``What Makes This Different from Ensemble-Based Approaches''
added in Section~I after the Contributions list.}

%----
\RC{The advantages and limitations of the proposed approach in relationship with
similar schemes is not clear.}

\RS{We have added a 7-column comparison table (Table~1 in the revised manuscript,
now at the end of Section~III) comparing Agentic SABRE against six representative
prior systems across: multi-agent design, uncertainty estimation, triage policy,
symbolic cues, explainability, concept-drift handling, and score-level augmentation.
On the limitations side, every section now opens with a 1--2 sentence statement of
its purpose and scope, and Section~XIII (Limitations and Future Work) now explicitly
opens by enumerating the three principal constraints on the current design.}

\CH{Comparison table (Table~1) added at the end of Section~III. Section-opening
sentences added to all 13 sections. Section~XIII opening expanded.}

%----
\RC{Please revise the structure of the paper. It is recommendable to add in each
section a couple of sentences that explain which is the purpose of the section.}

\RS{Done. Every section in the revised manuscript now opens with 1--2 sentences
stating its purpose within the overall narrative.}

\CH{Purpose sentences added to the opening of Sections~II through~XIII.}

\bigskip
\hrule
\bigskip

% -----------------------------------------------------------------------
%  REVIEWER 1
% -----------------------------------------------------------------------
\section{Reviewer \#1}

\subsection*{Weaknesses}

\RC{Weakness 1. The proposed approach is rather ensemble-based where components are
parallel or consequently connected DL models (CNN and GAN-based ones). The approach
itself looks more as an agent that provides decision of the sample rather than a
system of cooperative agents. Thus, the title is slightly confusing.}

\RS{We appreciate this important observation and address it directly in the new
``What Makes This Different'' subsection (Section~I). While the two CNNs are trained
independently and run in parallel at inference time, the system constitutes a genuine
multi-agent architecture in three respects that go beyond ensemble design: (1)~Each
agent specialises on a structurally distinct input modality---the Semantic Agent
processes PE-level embedding vectors ($\mathbb{R}^{384}$) while the Behavioural Agent
processes sliding-window I/O telemetry ($\mathbb{R}^{16}$). These are not bootstrap
resamples of the same feature space. (2)~The agents produce not only a posterior
mean probability but also an \emph{epistemic uncertainty} estimate via Monte Carlo
Dropout; this uncertainty is the primary mechanism by which the orchestrator decides
whether to act autonomously. A standard ensemble does not produce per-model
uncertainty signals of this kind. (3)~The Decision Agent (orchestrator) is a third
trained agent that operates over the score space of the first two, constituting a
hierarchical agent architecture. We have clarified this framing in the revised
Introduction.}

\CH{New subsection ``What Makes This Different from Ensemble-Based Approaches''
in Section~I; agent architecture clarified throughout Section~IV.}

%----
\RC{Weakness 2. The authors list the following contribution: ``lightweight symbolic
reasoning using ATT\&CK-aligned API ontologies and heuristic cues'' (4th item).
However, the ``ATT\&CK-aligned API ontologies'' are never referred to in the
approach description except for Sections II and III. This creates misleading
expectations. The reasoning is done based on heuristic threshold-based rules.}

\RS{The reviewer is correct. The original wording overstated the symbolic component.
We have rewritten the 4th contribution bullet to accurately describe what is
implemented: ATT\&CK-\emph{inspired} heuristic cues---specifically, API call
category indicators, entropy-threshold signals, and ransomware-relevant filesystem
access pattern flags---that are integrated into the feature extraction and triage
stages. We no longer claim a formalised ontology; instead we describe these as
lightweight symbolic enrichment signals grounded in ATT\&CK tactic categories.}

\CH{4th contribution bullet in Section~I rewritten. The revised text reads: ``We
incorporate \textbf{lightweight symbolic reasoning} via ATT\&CK-inspired heuristic
cues---including API call category indicators, entropy-threshold signals, and
ransomware-relevant filesystem access patterns---integrated into the feature
extraction and triage stages to enrich agent decisions with interpretable
threat-intelligence signals, forming a hybrid neuro-symbolic detection pipeline.''}

%----
\RC{Weakness 3a. Many statements are repeated in different sections: Sec.~III (the
last sentence), Sec.~V (the last paragraph in the beginning); Sec.~V~A (the second
paragraph); Sec.~V~C (the second paragraph).}

\RS{We have removed or reworked all identified repetitions. Specifically: the closing
sentence of Section~III (``This positions SABRE as a practical neuro-symbolic \ldots
rather than a conceptual demonstration'') has been rewritten to state, concretely,
which three gaps SABRE fills, rather than restating the system's positioning claim
that opens Section~V. The opening of Section~V now serves as the primary location for
this framing. Redundant restatements in Sections~V-A and V-C have been condensed.}

\CH{Section~III final paragraph rewritten. Redundant material in Sections~V-A
and V-C removed.}

%----
\RC{Weakness 3b. The presentation in terms of mathematical formulas is difficult
to follow: (i)~some formulas lack explanations to the notations used---what do
$p_{\text{beh}}, p_{\text{sem}}, \hat{p}$ stand for in Formula~(1)? What do
$s_{\text{beh}}, s_{\text{sem}}$ stand for in Formula~(1)? How does $\hat{p}$ in
Formula~(2) differ from $\hat{p}$ in Formula~(21)? (ii)~Formula~22 duplicates
Formula~3. (iii)~It is strange to have a reference to Formula~25 after Formula~10.}

\RS{We have addressed all three sub-points.}

\noindent\textbf{(i) Notation before Eq.~(1) and Eq.~(2).}
We added a prose paragraph before the score-vector equation that explicitly defines
$\bar{p}_{\text{sem}} \in [0,1]$ as the predictive mean of the Semantic Agent and
$\bar{p}_{\text{beh}} \in [0,1]$ as the predictive mean of the Behavioural Agent,
both obtained via MC Dropout. The notation $\mathbf{s}$ for the score vector is
introduced before it appears in any formula. The fused risk $\hat{p}$ in Eq.~(2)
is now explicitly described as the output of the MLP Decision Agent operating on
$\mathbf{s}$; a clarifying note explains that at deployment the full score vector
additionally incorporates per-agent uncertainties (introduced in Section~IX), so
the two uses of $\hat{p}$ are consistent---both refer to the MLP's output, with
the difference being whether uncertainties are included as inputs.

\noindent\textbf{(ii) Duplicate $\sigma_{\max}$ formula.}
The canonical definition of
$\sigma_{\max} = \max(\sigma_{\text{sem}}, \sigma_{\text{beh}})$ is now given once,
in Section~IV-C, with label \texttt{eq:sigmamax}. The occurrence in Section~IX
that duplicated it has been replaced by a prose reference to that equation.

\noindent\textbf{(iii) Forward reference to $J(\tau,\kappa)$.}
The scalarised objective $J(\tau,\kappa) = \mathrm{bal\_acc}(\tau,\kappa) -
\lambda\,\mathrm{EscRate}(\tau,\kappa)$ has been moved from Section~IX to
Section~VI (Policy Analysis), where it is first discussed. Section~IX now
references Eq.~(12) in Section~VI, eliminating the forward
reference.

\CH{Variable definitions added before score-vector formula in Section~IV-B.
Duplicate $\sigma_{\max}$ equation removed from Section~IX. $J(\tau,\kappa)$
definition relocated from Section~IX to Section~VI.}

%----
\RC{Weakness 4. How are thresholds $\tau$, $\tau_{\text{high}}$ ordered? How are
thresholds $\kappa$, $\kappa_{\text{low}}$ ordered?}

\RS{We have added an explicit ordering constraint as a numbered equation in
Section~V-A:
$\tau \le \tau_{\text{high}}$ and $\kappa_{\text{low}} \le \kappa$.
The prose explanation now states clearly that $\tau_{\text{high}}$ demands a higher
fused risk for autonomous containment than $\tau$ does for mere escalation, and that
$\kappa_{\text{low}}$ imposes a tighter (lower) uncertainty ceiling than $\kappa$.}

\CH{Equation with ordering constraints $\tau \le \tau_{\text{high}}$,
$\kappa_{\text{low}} \le \kappa$ added in Section~V-A with accompanying explanation.}

%----
\RC{Weakness 5. Figure~5 requires additional labels to axis and dotted plots.}

\RS{We acknowledge this. In the revised manuscript, the caption of Figure~5 has been
expanded to explicitly name each threshold line ($\tau$, $\tau_{\text{high}}$,
$\kappa$, $\kappa_{\text{low}}$) and their roles. The figure itself will be
regenerated with annotated threshold labels before the camera-ready version, as the
figure generation scripts are maintained in the accompanying codebase.}

\CH{Figure~5 caption expanded to annotate all four threshold lines.
Figure regeneration with embedded labels is planned for the camera-ready version.}

%----
\RC{Weakness 6. Why are datasets included in the approach description? Is the
approach dataset-dependent?}

\RS{The approach is \emph{not} dataset-dependent. Section~VIII (Experimental Setup)
now opens with an explicit statement to this effect, and the datasets subsection
now clarifies that RDset and RanSMAP are used to \emph{instantiate} the two
complementary modalities (static PE semantics and runtime I/O telemetry) that the
approach assumes in the abstract. The architecture, triage policy, and all algorithms
are specified entirely in terms of generic inputs $D_{\text{sem}}$ and $D_{\text{beh}}$
in Algorithms~1--4 and are fully dataset-agnostic.}

\CH{Section~VIII now opens with a statement clarifying that the approach is
dataset-independent and that the specific datasets are used only to instantiate
the two generic complementary modalities.}

%----
\RC{Weakness 7. Do datasets contain the same malware samples? How are features from
different datasets fused? What labels are used to construct the dataset for the
GAN-based model?}

\RS{We have added a dedicated paragraph titled \textbf{``Cross-source fusion
rationale''} in Section~VIII-A addressing all three questions. In brief:
(1)~RDset and RanSMAP contain \emph{no overlapping malware samples}---RDset
consists of PE file metadata from Windows executables, while RanSMAP contains
runtime I/O traces recorded in instrumented hypervisors. (2)~Raw features are
\emph{never} merged: each agent trains on its own dataset, and fusion occurs only
at the \emph{score level} (the two posterior mean probabilities are concatenated
into $\mathbf{s}$). (3)~CTGAN is conditioned on binary labels
$y_i \in \{0,1\}$ (benign / ransomware) from the training set; no label
transformations are applied.}

\CH{``Cross-source fusion rationale'' paragraph and ``CTGAN labels'' paragraph
added in Section~VIII-A.}

%----
\RC{Weakness 8. What is the ``predefined low risk threshold $\theta$'' (p.~9, first
line, Sec.~IX-B)? How is it defined? What is its value?}

\RS{We have added a new subsection (Section~IX-A, ``Counterfactual Threshold
$\theta$ and Behavioural Decision Boundary'') that defines $\theta$ explicitly:
it is the maximum posterior probability at which the agent's output is classified
as ``benign'' for the purpose of counterfactual generation. We set $\theta = 0.49$,
just below the standard 0.5 decision boundary, to enforce a strict label flip.
The subsection also reports that the 100\% flip success rate is stable across
$\theta \in [0.40, 0.49]$, confirming low sensitivity to this choice.}

\CH{New Section~IX-A ``Counterfactual Threshold $\theta$ and Behavioural Decision
Boundary'' added, defining $\theta = 0.49$ with sensitivity note.}

\bigskip
\hrule
\bigskip

% -----------------------------------------------------------------------
%  REVIEWER 2
% -----------------------------------------------------------------------
\section{Reviewer \#2}

\RC{Point 1. The manuscript should more clearly distinguish its contribution from
prior ensemble-based ransomware detection systems. A concise comparison paragraph
highlighting what is fundamentally new would improve positioning.}

\RS{Addressed by the new ``What Makes This Different from Ensemble-Based Approaches''
subsection in Section~I and the 7-column comparison table (Table~1) in Section~III.
The subsection explains the three architectural properties that make SABRE a genuine
multi-agent system rather than an ensemble: heterogeneous input spaces, per-agent
uncertainty estimation used as an operational control signal, and a formal executable
triage policy.}

\CH{New subsection in Section~I; comparison Table~1 in Section~III.}

%----
\RC{Point 2. The semantic agent achieves AUC$=1.0000$ on RDset. The manuscript should
include a clearer discussion on generalisation risks (family-wise split, temporal
split, structural similarity bias).}

\RS{We have added a paragraph titled \textbf{``Generalisation caveats for
AUC$=1.0$''} in Section~IX immediately following the AUC result. This paragraph
explicitly acknowledges: (1) that RDset does not employ a family-wise or temporal
split, so structural similarity bias is possible; (2) that the perfect AUC reflects
the highly structured nature of PE-level metadata in an \emph{intra-family}
evaluation setting; and (3) that the saturated semantic regime is deliberately used
as a stress-test for the fusion and triage layers, not as a claim of general
dominance. Generalisation to temporally displaced or family-disjoint corpora is
identified as future work.}

\CH{``Generalisation caveats for AUC$=1.0$'' paragraph added in Section~IX.}

%----
\RC{Point 3. The behavioural agent reports high accuracy but relatively modest AUC.
This discrepancy should be more clearly explained, particularly regarding class
imbalance, threshold choice, and probability calibration.}

\RS{We have expanded the Behavioural Agent result paragraph in Section~IX with a
detailed explanation of this discrepancy. In brief: the high accuracy is driven by
the class prior (most samples are benign, so a low-probability classifier is
accurate by default); AUC measures rank separability regardless of prior, and the
modest AUC reflects the genuine overlap between malicious and benign I/O distributions
in a polymorphic telemetry stream. SMOTE re-balancing improves training stability but
does not resolve the distributional overlap. Crucially, the Behavioural Agent's role
is framed as a \emph{modulator of $\sigma_{\max}$} rather than a standalone detector:
its diffuse posteriors raise uncertainty for borderline cases, routing them to human
analysts rather than enabling overconfident autonomous containment.}

\CH{Expanded explanation of the accuracy/AUC discrepancy added in Section~IX
following the Behavioural Agent result.}

%----
\RC{Point 4. A short discussion under Section~IV comparing Monte Carlo Dropout with
alternative uncertainty approaches (deep ensembles, calibration stability under
distribution shift) would strengthen deployment credibility.}

\RS{We have added a paragraph titled \textbf{``Choice of uncertainty estimator''}
in Section~IV-C (Uncertainty Aggregation). This paragraph explains why MC Dropout
was chosen over deep ensembles and Bayesian neural networks: no architectural
changes, lower inference cost ($T$ forward passes versus $M$ independently trained
models), and compatibility with the existing dropout regularisation. The paragraph
also honestly acknowledges that deep ensembles can produce better-calibrated
uncertainty under strong distribution shift, and notes this as a limitation addressed
in Section~XIII.}

\CH{``Choice of uncertainty estimator'' paragraph added in Section~IV-C.}

%----
\RC{Point 5. The safety-optimal threshold $\tau = 0.999999$ appears extremely strict.
The manuscript should clarify whether this value is dataset-specific and include a
brief sensitivity or stability analysis.}

\RS{We have added a paragraph titled \textbf{``Threshold sensitivity and dataset
specificity''} in Section~IX. This paragraph makes three points: (1)~$\tau$ is
dataset-specific---the near-unity value reflects the bimodal ($\{0,1\}$-clustering)
nature of the fused probabilities on RDset, and a lower $\tau$ would be appropriate
for datasets with more dispersed posteriors; (2)~performance is robust to the exact
value: both the safety-optimal and low-escalation policies yield near-identical
accuracy (0.99999 on RDset, 0.9406 vs.\ 0.9405 on RanSMAP), showing headline metrics
are stable; (3)~the triage geometry is inherently stable near $\tau$ because the
boundary neighbourhood in $(\hat{p},\sigma_{\max})$ space is sparsely populated
(argued in Section~VI).}

\CH{``Threshold sensitivity and dataset specificity'' paragraph added in Section~IX.}

%----
\RC{Point 6. The adversarial cost analysis is strong. However, counterfactual norms
operate in continuous space. A short clarification of how these costs translate to
real executable constraints (PE format restrictions, runtime feasibility) would
further enhance practical cybersecurity relevance.}

\RS{We have added a paragraph titled \textbf{``Translation of counterfactual costs
to real-world executable constraints''} in Section~X, immediately before the Failure
Modes discussion. For the Behavioural Agent, we explain that reducing LBA/GPA
variance requires altering the ransomware's I/O access pattern, which conflicts
directly with its encryption strategy (ATT\&CK T1486: Data Encrypted for Impact),
creating a functional trade-off. For the Semantic Agent, we explain that perturbing
PE structural features requires maintaining PE header validity (Magic bytes,
SizeOfCode consistency), bounding the feasible deformation to a subset of the
continuous neighbourhood $\mathcal{N}_\epsilon(\mathbf{z})$. The reported
$\mathbb{E}[\mathcal{C}_{\text{sem}}] = 3.27$ is therefore a \emph{lower bound};
the true adversarial cost under executable constraints is higher.}

\CH{``Translation of counterfactual costs to real-world executable constraints''
paragraph added in Section~X before the Failure Modes subsection.}

\bigskip
\hrule
\bigskip

% -----------------------------------------------------------------------
%  REVIEWER 4
% -----------------------------------------------------------------------
\section{Reviewer \#4}

\RC{Point 1. The claimed innovation is somewhat overstated. The authors are encouraged
to clarify the substantive technical advances over existing uncertainty-aware
detection, multi-agent detection, and neuro-symbolic approaches.}

\RS{We have restructured the Introduction to foreground the three substantive
novelties more precisely. The new ``What Makes This Different'' subsection explains
that: prior uncertainty-aware detectors use uncertainty as a \emph{diagnostic
output} (Kabuye et al.~\cite{kabuye2025explainable}), whereas SABRE uses it as an
\emph{operational control signal} that formally constrains autonomous action;
prior multi-agent detectors fuse votes or feature concatenations (Soltani et al.,
2023), whereas SABRE fuses heterogeneous modalities at the score level under joint
uncertainty governance; and prior neuro-symbolic systems embed rules at the
architecture level, whereas SABRE embeds them as interpretable triage thresholds
that can be updated without model retraining. The comparison Table~1 quantifies
these distinctions across six reference systems.}

\CH{New subsection in Section~I; comparison Table~1 in Section~III; contributions
list in Section~I revised for precision.}

%----
\RC{Point 2. The dataset setting is evidently heterogeneous. The rationale for
cross-source fusion, the comparability of these datasets, and the real-world
deployment scenario should be further clarified.}

\RS{The ``Cross-source fusion rationale'' paragraph in Section~VIII-A directly
addresses this. In brief: RDset and RanSMAP are not comparable in feature space---they
are complementary observation planes of the same threat category (static PE artefacts
vs.\ runtime I/O telemetry). Fusion at the raw feature level would be meaningless;
fusion at the score level is motivated by the realistic SOC scenario in which a
static sandbox scanner and a runtime endpoint monitor report simultaneously, and a
security orchestrator must act on both signals. This mirrors real-world deployments
where SIEM platforms ingest alerts from heterogeneous sensors. The Experimental Setup
section now opens with this scenario framing.}

\CH{``Cross-source fusion rationale'' paragraph in Section~VIII-A; Section~VIII
opening paragraph revised to frame the real-world deployment scenario.}

%----
\RC{Point 3. The threshold strategy relies heavily on parameters $\tau$, $\kappa$,
$\tau_{\text{high}}$, $\kappa_{\text{low}}$. The parameter selection process,
sensitivity analysis, and cross-dataset stability are insufficient.}

\RS{We have addressed this in three ways. First, the policy objective
$J(\tau,\kappa)$ is now defined in Section~VI (Policy Analysis) with a full
description of the optimisation procedure (grid search over the non-smooth landscape
with $\lambda = 0.5$). Second, the ``Threshold sensitivity and dataset specificity''
paragraph in Section~IX explains why $\tau=0.999999$ is dataset-specific and
demonstrates stability: both safety-optimal and low-escalation policies achieve
near-identical accuracy across both datasets, confirming robustness to small
perturbations in the threshold pair. Third, Section~VI explains the geometric
stability argument: the boundary neighbourhood in $(\hat{p},\sigma_{\max})$ space
is sparsely populated due to the monotonic fusion structure, meaning small threshold
shifts move the boundary without reclassifying many samples. We acknowledge that
a rigorous cross-dataset stability evaluation (e.g.\ on held-out ransomware families)
remains an important future direction.}

\CH{$J(\tau,\kappa)$ definition moved to Section~VI; threshold sensitivity paragraph
added in Section~IX; geometric stability argument in Section~VI.}

%----
\RC{Point 4. The manuscript repeatedly emphasises calibration, interpretability, and
robustness. However, the current evidence is mainly based on qualitative analysis or
limited case examples. More systematic statistical tests and comparative experiments
are recommended.}

\RS{We acknowledge this limitation. The existing evidence is structured across
three complementary quantitative channels: (i)~reliability diagrams confirming
$\mathbb{P}(y=1 \mid \hat{p}=q) \approx q$ (Figure~3); (ii)~ablation tables
(Tables~VIII and IX) that quantify the impact of removing each component;
(iii)~$N=50$ counterfactual experiments with distributional statistics
(Tables~VII and~III). We agree that formal statistical tests (e.g.\ McNemar's
test on the confusion matrices of the full model versus each ablated variant,
or bootstrap confidence intervals on AUC differences) would further strengthen
validity claims. We have noted this explicitly as a future-work item in
Section~XIII and have added language clarifying the nature of the evidence
(empirical rather than hypothesis-test-level) at the start of Section~IX.}

\CH{Section~IX opening paragraph clarified regarding the nature of evidence;
Section~XIII updated to list formal statistical validation as a future direction.}

%----
\RC{Point 5. The adversarial robustness analysis is mainly based on counterfactual
perturbation norms and theoretical discussion. Standard adversarial attack
experiments are lacking.}

\RS{We agree that standard adversarial attack experiments (e.g.\ FGSM or PGD on
the neural agents) would strengthen robustness claims. Executing such attacks on
the behavioural agent requires a differentiable model of the I/O simulator, and
on the semantic agent requires a generative model of PE binaries that satisfies
loader constraints---both are substantial engineering efforts outside the scope of
the current paper. The counterfactual analysis serves as a principled surrogate:
it computes the \emph{minimal-norm perturbation to cross the decision boundary},
which is equivalent to the white-box $\ell_2$-norm adversarial attack under a
continuous relaxation. The ``Translation to real-world executable constraints''
paragraph (Section~X) now grounds these norms in concrete PE and I/O constraints.
We have added a sentence in Section~XIII acknowledging standard adversarial attack
evaluation as a priority for future work.}

\CH{Section~X adversarial cost paragraph expanded with real-world grounding.
Section~XIII updated to identify standard adversarial evaluation as future work.}

\bigskip
\hrule
\bigskip

% -----------------------------------------------------------------------
%  REVIEWER 5
% -----------------------------------------------------------------------
\section{Reviewer \#5}

\RC{The experimental results section would benefit from including confidence
intervals or variance metrics to further support the uncertainty-aware claims.}

\RS{Monte Carlo Dropout itself produces per-sample uncertainty estimates
$\sigma_{\text{sem}}$ and $\sigma_{\text{beh}}$ that can be interpreted as
predictive standard deviations; these are summarised in Table~VII for the
counterfactual experiments. For aggregate performance metrics (AUC, accuracy),
formal bootstrap confidence intervals require repeated held-out evaluation which
is computationally intensive at the scale of RanSMAP (11,820 execution traces).
We have added a clarifying note in Section~IX stating that the MC Dropout
uncertainties constitute the uncertainty evidence at the per-sample level, and
that bootstrap CI computation for aggregate metrics is planned for future work.
We also note that the ablation comparisons in Tables~VIII and IX provide a form
of model-level variance assessment across architecturally distinct configurations.}

\CH{Clarifying note on confidence intervals added in Section~IX.
Section~XIII updated with formal CI computation as a future direction.}

%----
\RC{A brief discussion on computational complexity and runtime performance would
improve the practical relevance of the work.}

\RS{We have substantially expanded the subsection titled \textbf{``Computational
Complexity and Runtime Performance''} in Section~XII (Discussion) to include a
formal Big-O analysis of every pipeline component, structured in six paragraphs
and four numbered equations.
Specifically:
\begin{itemize}[leftmargin=1.5em,itemsep=2pt]
  \item \textbf{CNN agent training} (Eq.~(55)):
        $\mathcal{O}(E \cdot N \cdot L \cdot K \cdot C^2)$ per agent, linear in
        both sample count $N$ and epoch count $E$, with explicit constants
        $L=3$, $K=3$, $C=64$, and modality-specific input dimensions
        ($C_0=384$ semantic; $C_0=16$ behavioural).
  \item \textbf{MC Dropout score generation}:
        $\mathcal{O}(2 T N L K C^2)$ for both agents combined, with $T=30$
        stochastic passes; empirical wall-clock: $\approx 60$--$150$\,ms per
        agent on CPU.
  \item \textbf{CTGAN augmentation} (Eq.~(56)):
        $\mathcal{O}(E_{\mathrm{CTGAN}} \cdot N_s \cdot
        [\sum_i g_i g_{i+1} + \sum_j d_j d_{j+1}])$; the score matrix is
        $N_s \times 2$---orders of magnitude smaller than the raw feature
        matrices---so 200 epochs fit in $\approx 5$--$10$\,min on CPU.
  \item \textbf{Decision Agent MLP} (Eq.~(57)):
        $\mathcal{O}(E_{\mathrm{MLP}} \cdot N_{\mathrm{aug}} \cdot h^2)$,
        negligible relative to CTGAN given the small input dimension ($d=4$).
  \item \textbf{Bayesian TPE optimisation} (Eq.~(58)):
        $\mathcal{O}(T_{\mathrm{TPE}} \cdot (n_{\mathrm{obs}} \log n_{\mathrm{obs}}
        + N_{\mathrm{val}}))$, converging within minutes.
  \item \textbf{Per-sample inference}: dominated by MC Dropout at
        $\mathcal{O}(T \cdot L \cdot K \cdot C^2)$; total wall-clock
        $\approx 0.15$--$0.35$\,s CPU, $<10$\,ms GPU.
  \item \textbf{Total training complexity} (Eq.~(59)):
        a four-term closed-form expression covering all components, revealing
        that each CNN agent and the triage thresholds can be updated independently
        without retraining the rest of the pipeline.
\end{itemize}}

\CH{``Computational Complexity and Runtime Performance'' subsection in
Section~XII substantially expanded with formal Big-O analysis of all pipeline
components (four new numbered equations: (55),
(56), (57), (58),
(59)) and six structured paragraphs.}

%----
\RC{The related work section may be strengthened by including a brief comparison
table summarising key differences between existing approaches and the proposed
framework.}

\RS{Done. Table~1 in the revised Section~III compares Agentic SABRE against six
representative systems across seven design dimensions: multi-agent architecture,
uncertainty estimation, operational triage policy, symbolic cues, explainability,
concept-drift handling, and score-level augmentation. The table uses a three-level
encoding (\checkmark~= present; $\circ$~= partial; $\times$~= absent) for each
dimension.}

\CH{Comparison table (Table~1) added at the end of Section~III.}

%----
\RC{The manuscript is generally well-written; however, a few grammatical
inconsistencies and typographical errors should be carefully reviewed.}

\RS{The entire manuscript has been re-read and edited for grammatical consistency
and typographical accuracy. Key corrections include: uniform use of the en-dash
vs.\ hyphen in compound adjectives; removal of doubled words and inconsistent
capitalisation in section headings; and standardised notation for all mathematical
symbols (e.g.\ consistent use of $\bar{p}$ vs.\ $\hat{p}$ and $\sigma$ vs.\
$\sigma_{\max}$ across sections).}

\CH{Manuscript-wide grammar and typographical review completed; corrections
applied throughout.}

\bigskip
\hrule
\bigskip

% -----------------------------------------------------------------------
%  CLOSING
% -----------------------------------------------------------------------
\section*{Closing Remark}

We believe the revised manuscript is substantially stronger in its technical
precision, structural clarity, and contextual positioning. We are grateful to all
reviewers for the depth of their engagement. We remain available to address any
remaining concerns the editor or reviewers may have.

\bigskip
\noindent\textit{The Authors}

\end{document}
